\documentclass[12pt]{article}
\usepackage[utf8]{inputenc}
\usepackage[T1]{fontenc}
\usepackage{hyperref}
\title{Yuchi Language Resources for K-12 Curriculum Development}
\author{Compiled by Claude Sonnet 5.5 (claude-sonnet-5-5) for Andr\'as Kornai}
\date{}
\begin{document}
\maketitle

\section{Introduction}
Yuchi (Euchee; ISO 639-3 yuc) is a language isolate of Oklahoma with very few elderly first-language speakers. Older sources give conflicting speaker counts (about 12 elders, or 5 by 2010). This report rests on a small number of web search results. Little was found, and nothing here is a ready-made K-12 curriculum.

\section{Primary Educational Resources}
The Euchee (Yuchi) Language Project in Sapulpa, Oklahoma, runs an after-school program for children aged 3 to 18, a Master-Apprentice program, and a weekly community class. Coverage of the project is at \url{https://news.harvard.edu/gazette/story/2017/11/yuchi-language-project-seeks-to-revitalize-indigenous-identity-through-common-tongue/} and \url{https://atomicscribe.com/we-are-still-here-the-yuchi-language-project/}. No textbooks or primers were found in the search results.

\section{Linguistic and Reference Materials}
An NSF grant (BCS-0345694) planned a dictionary of about 6,000 entries plus a text and sound database: \url{https://grantome.com/grant/NSF/BCS-0345694}. The search results do not show that the dictionary was published. The Wikipedia page is at \url{https://wikipedia.com/wiki/ISO_639:yuc}.

\section{Religious and Community Resources}
No scripture translations were found. A community-oriented article on culture and language is at \url{https://indianyouth.org/2014/07/11/}.

\section{Digital and Audio Resources}
The Endangered Languages Project lists Yuchi resources, samples and a guide: \url{https://www.endangeredlanguages.com/node/37636/resources}, \url{https://endangeredlanguages.com/lang/1002/samples/11604?hl=en}, \url{https://endangeredlanguages.com/lang/1002/guide}.

\section{Curriculum Development Context}
Immersion-style community programs exist, but the language is not known to be taught in schools. A curriculum would have to be built together with the Euchee Language Project and the few remaining speakers, from the documentation above.

\end{document}